\documentclass[11pt]{amsart}
\usepackage[margin=27mm]{geometry}
\usepackage{amsmath,amssymb,amsthm}
\usepackage{array,booktabs}
\usepackage[obeyspaces,spaces]{url}
\usepackage[colorlinks=true,linkcolor=blue,citecolor=blue,urlcolor=blue]{hyperref}
% The author's macros.
%
% some calligraphy
%
\newcommand\cA{{\mathcal A}}
\newcommand\cB{{\mathcal B}}
\newcommand\cC{{\mathcal C}}
\newcommand\cD{{\mathcal D}}
\newcommand\cE{{\mathcal E}}
\newcommand\cF{{\mathcal F}}
\newcommand\cG{{\mathcal G}}
\newcommand\cH{{\mathcal H}}
\newcommand\cI{{\mathcal I}}
\newcommand\cJ{{\mathcal J}}
\newcommand\cK{{\mathcal K}}
\newcommand\cL{{\mathcal L}}
\newcommand\cM{{\mathcal M}}
\newcommand\cN{{\mathcal N}}
\newcommand\cP{{\mathcal P}}
\newcommand\cQ{{\mathcal Q}}
\newcommand\cR{{\mathcal R}}
\newcommand\cS{{\mathcal S}}
\newcommand\cT{{\mathcal T}}
\newcommand\cU{{\mathcal U}}
\newcommand\cV{{\mathcal V}}

%
% fields and rings (and a semigroup)
%
\newcommand\cx{{\mathbb C}}% complexes
\newcommand\fld{{\mathbb F}}
\newcommand\flde{{\mathbb E}}
\newcommand\ints{{\mathbb Z}}
\newcommand\nn{{\mathbb N}}%non-negative integers
\newcommand\re{{\mathbb R}}%reals
\newcommand\rats{{\mathbb Q}}

\newcommand\Zu{{\mathbf u}}
\newcommand\Zv{{\mathbf v}}
\newcommand\Zw{{\mathbf w}}
\newcommand\Zy{{\mathbf y}}
\newcommand\Zz{{\mathbf z}}
\newcommand\Zx{{\mathbf x}}
\newcommand\Zj{{\mathbf j}}


%
% the really useful stuff
%
%\newcommand\opk[1]{\mathop{\hbox{\rm #1}}\nolimits} %% old version
\newcommand\opk[1]{\mathop{\mathrm{#1}}\nolimits}
\newcommand\comp[1]{{\mkern2mu\overline{\mkern-2mu#1}}}
\newcommand\diff{\mathbin{\mkern-1.5mu\setminus\mkern-1.5mu}}% for \setminus
\newcommand\gbin[2]{\genfrac{[}{]}{0pt}{}{#1}{#2}}
\newcommand\inv{^{-1}}
\newcommand\res{\mathbin{\mkern-2.0mu\restriction\mkern-2.0mu}}% ch2
\newcommand\sbs{\subseteq}
\newcommand\sps{\supseteq}
\newcommand\seq[3]{#1_{#2},\ldots,#1_{#3}}
\newcommand\supp{\opk{supp}}
\newcommand\im{\opk{im}}
\newcommand\row{\opk{row}}
\newcommand\pmat[1]{\begin{pmatrix} #1 \end{pmatrix}}
\newcommand\cprod{\mathbin{\square}}
%
% matrix theory
%
\newcommand\ip[2]{\langle#1,#2\rangle}
\newcommand\one{{\bf1}}
\newcommand\rk{\opk{rk}}
\newcommand\tr{\opk{tr}}
\newcommand\col{\opk{col}}
\newcommand\mat[3]{\mathrm{Mat}_{#1\times #2}(#3)}

%
% some group theory
%
\newcommand\aut[1]{{\rm Aut}(#1)}
\newcommand\fx[1]{{\rm fix}(#1)}% ch2
\newcommand\grp[1]{\langle #1\rangle}
\newcommand\nrml{\vartriangleleft}
\newcommand\nrmleq{\trianglelefteq}
\newcommand\alt[1]{\opk{Alt}(#1)}
\DeclareMathOperator{\Sym}{Sym}
% End of the author's macros.
\newtheorem{theorem}{Theorem}[section]
\newtheorem{lemma}[theorem]{Lemma}
\newtheorem{proposition}[theorem]{Proposition}
\newtheorem{corollary}[theorem]{Corollary}
\theoremstyle{definition}
\newtheorem{definition}[theorem]{Definition}
\newtheorem{remark}[theorem]{Remark}
\DeclareRobustCommand{\ml}[1]{\ifmmode\text{\normalfont\small\ttfamily\detokenize{#1}}\else{\normalfont\small\nolinkurl{#1}}\fi}
\DeclareUnicodeCharacter{03BC}{\ensuremath{\mu}}
\DeclareUnicodeCharacter{2113}{\ensuremath{\ell}}
\DeclareUnicodeCharacter{1D9C}{\ensuremath{^{\mathrm c}}}
\DeclareUnicodeCharacter{2243}{\ensuremath{\simeq}}
\DeclareUnicodeCharacter{2206}{\ensuremath{\triangle}}
\DeclareUnicodeCharacter{00D7}{\ensuremath{\times}}
\DeclareUnicodeCharacter{2295}{\ensuremath{\oplus}}
\DeclareUnicodeCharacter{0393}{\ensuremath{\Gamma}}
\DeclareUnicodeCharacter{21A6}{\ensuremath{\mapsto}}
\DeclareUnicodeCharacter{2229}{\ensuremath{\cap}}
\DeclareUnicodeCharacter{2086}{\ensuremath{_6}}
\DeclareUnicodeCharacter{2087}{\ensuremath{_7}}
\newcommand{\mlu}[1]{{\normalfont\small\ttfamily #1}}
\newcommand{\nb}{N}
\newcommand{\sd}{\mathbin{\triangle}}
\emergencystretch=2em
\title[The Clebsch and Schl\"afli graphs]{The Clebsch and Schl\"afli graphs}
\date{28 September 2026}

\begin{document}
\raggedbottom
\begin{abstract}
A strongly regular graph with parameters $(16,5,0,2)$ is isomorphic to the folded 5-cube, and so
one with parameters $(16,10,6,6)$ is isomorphic to the Clebsch graph; a strongly regular graph
with parameters $(27,16,10,8)$ is isomorphic to the Schl\"afli graph. These are prerequisites SE2
and SE3 of \ml{cgss/prerequisites.md}, used in step $\langle1\rangle11$ of Remark 4.14 of
[CGSS]. We give a direct proof of the first, reduce the second to the uniqueness assertion of
Lemma 2.10 of [CGSS] as repaired in the blueprint, check both with SageMath, and describe the
shape of the Lean statements.
\end{abstract}
\maketitle

\section{Statement}\label{sec:statement}

All graphs are simple and finite. For a vertex $v$ of a graph $X$ we will write $\nb(v)$ for its
neighbourhood, $u\sim w$ when $u$ and $w$ are adjacent, and $\comp X$ for the complement of $X$.
A graph $X$ on $n$ vertices is \emph{strongly regular with parameters $(n,k,a,c)$} if it is
$k$-regular, any two adjacent vertices have exactly $a$ common neighbours, and any two distinct
nonadjacent vertices have exactly $c$ common neighbours. This is Mathlib's
\ml{SimpleGraph.IsSRGWith n k a c}, with no further condition; in particular $X$ may be complete
or empty. For sets $s,t$ we will write $s\sd t$ for the symmetric
difference, and for a set $\Omega$ we will write $\binom{\Omega}{2}$ for the set of its 2-subsets.

\begin{definition}\label{def:clebsch}
Let $\Omega_5=\{1,\dots,5\}$, and let $V_{16}$ be the set of the subsets of $\Omega_5$ of even
size; thus $V_{16}$ consists of $\emptyset$, the ten 2-subsets and the five 4-subsets, and
$\lvert V_{16}\rvert=16$. The \emph{folded 5-cube} $\square_5$ is the graph on $V_{16}$ in which
$s\sim t$ if and only if $\lvert s\sd t\rvert=4$. The \emph{Clebsch graph} is its complement
$\comp{\square_5}$.
\end{definition}

For distinct $s,t\in V_{16}$ the size $\lvert s\sd t\rvert=\lvert s\rvert+\lvert t\rvert-2\lvert
s\cap t\rvert$ is even, positive and at most 5, so it is 2 or 4. Thus $s$ and $t$ are adjacent
in the Clebsch graph if and only if $\lvert s\sd t\rvert=2$, and the Clebsch graph is the halved
5-cube. The folded 5-cube is the 5-cube with antipodal vertices identified: each pair $\{x,\Omega_5\diff
x\}$ contains exactly one set of even size, and $x\mapsto x\cup\{i\}$ or $x\diff\{i\}$ becomes $s\mapsto
\Omega_5\diff(s\sd\{i\})$, at distance 4 from $s$. Neither description is needed below; both are
checked in Section~\ref{sec:sage}. We follow [CGSS] Example 4.8 and \ml{cgss/notation.md} \S6 in
giving the name to the 10-regular graph; SageMath's \ml{graphs.ClebschGraph()} is the 5-regular
one, our $\square_5$. [CGSS] Example 4.8 gives the Clebsch graph by the 16 vectors
$\mathbf e_i+\mathbf e_5$, $\mathbf e_i+\mathbf e_6$, $\mathbf c-\mathbf e_i-\mathbf e_7$ and
$\mathbf c-\mathbf e_i-\mathbf e_8$ ($1\le i\le4$) of $\mathrm E_8$, with $\mathbf
c=\frac12\sum_{m=1}^8\mathbf e_m$, two of them adjacent when their inner product is 1; this graph
is isomorphic to $\comp{\square_5}$ (Section~\ref{sec:sage}), and in Lean the two are different
graphs joined by an explicit isomorphism (Section~\ref{sec:lean}).

\begin{definition}\label{def:schlafli}
Let $\Omega_6=\{1,\dots,6\}$, and let $V_{27}$ be the set of the 27 symbols $a_i$, $b_i$
($i\in\Omega_6$) and $c_{ij}=c_{ji}$ ($\{i,j\}\in\binom{\Omega_6}{2}$). Let $Q$ be the graph on
$V_{27}$ with the following edges, and no others:
\begin{enumerate}
\item $a_i\sim b_j$ for $i\ne j$;
\item $a_i\sim c_{jk}$ and $b_i\sim c_{jk}$ for $i\in\{j,k\}$;
\item $c_{ij}\sim c_{kl}$ for $\{i,j\}\cap\{k,l\}=\emptyset$.
\end{enumerate}
The \emph{Schl\"afli graph} is its complement $\comp Q$. Thus in $\comp Q$ we have $a_i\sim a_j$
and $b_i\sim b_j$ for $i\ne j$, $a_i\sim b_i$, $a_i\sim c_{jk}$ and $b_i\sim c_{jk}$ for
$i\notin\{j,k\}$, and $c_{ij}\sim c_{kl}$ when $\lvert\{i,j\}\cap\{k,l\}\rvert=1$.
\end{definition}

This is Schl\"afli's double six: $a_i$, $b_i$ and $c_{ij}$ are the 27 lines on a smooth cubic
surface, $Q$ is the graph of meeting lines, and the Schl\"afli graph is the graph of skew lines.
The triangles of $Q$ are the 45 tritangent planes, and $Q$ is the collinearity graph of the
generalized quadrangle of order $(2,4)$. We will not use this. [BCN] 3.11 (ii) gives the same graph in
coordinates: $a_i=\mathbf e_i+\mathbf e_7$, $b_i=\mathbf e_i+\mathbf e_8$ and $c_{ij}=\mathbf
c-\mathbf e_i-\mathbf e_j$ in $\mathrm E_8$, with $\mathbf c=\frac12\sum_{m=1}^8\mathbf e_m$, and
two vertices adjacent in $\comp Q$ when their inner product is 1 (it is 0 otherwise); this too is
checked in Section~\ref{sec:sage}.

[CGSS] Example 4.6, with $\mathbf c$ in place of the misprinted $\frac13\sum_k\mathbf e_k$
(\ml{cgss/errata.md}, Error 2), represents the Schl\"afli graph by other vectors of $\mathrm
E_8$: $\mathbf x_{ij}=\mathbf e_i+\mathbf e_j$ for $1\le i<j\le6$, and $\mathbf y_i=\mathbf
c-\mathbf e_i-\mathbf e_7$ and $\mathbf z_i=\mathbf c-\mathbf e_i-\mathbf e_8$ for $1\le i\le6$,
two of them adjacent when their inner product is 1. The bijection
\[
a_i\mapsto\mathbf y_i=\mathbf c-\mathbf e_i-\mathbf e_7,\qquad
b_i\mapsto\mathbf z_i=\mathbf c-\mathbf e_i-\mathbf e_8,\qquad
c_{ij}\mapsto\mathbf x_{ij}=\mathbf e_i+\mathbf e_j
\]
is an isomorphism from $\comp Q$ onto the graph of these vectors: for $i\ne k$ we have $(\mathbf
y_i,\mathbf y_k)=(\mathbf z_i,\mathbf z_k)=1$; $(\mathbf y_i,\mathbf z_k)$ is 1 if $i=k$ and 0
otherwise; $(\mathbf x_{ij},\mathbf y_k)=(\mathbf x_{ij},\mathbf z_k)$ is 1 if $k\notin\{i,j\}$
and 0 otherwise; and $(\mathbf x_{ij},\mathbf x_{kl})=\lvert\{i,j\}\cap\{k,l\}\rvert$ (the
blueprint, Example 4.6, step \ml{s46ip}). These are the adjacencies of $\comp Q$ listed in
Definition~\ref{def:schlafli}. The vectors of [BCN] and of [CGSS] are different, but under the
two labellings they have the same Gram matrix. The correspondence is checked in
Section~\ref{sec:sage}, and it is the certificate isomorphism of Section~\ref{sec:lean}.

\emph{Why this model.} Of the three usual descriptions (the 27 lines, the 27 vectors of
$\mathrm E_8$ above, and a quadric or Hermitian model over a finite field), the first is the
easiest to state in Lean: its vertex type is a sum of two finite types, its adjacency is
decidable by cases on four set conditions, and it needs no inner products and no field.
Its vertices correspond one to one to the vectors of the second and to those of [CGSS]
Example 4.6, so that the representation of the Schl\"afli graph in $\mathrm E_8$ becomes a check
on explicit vectors of the form of step $\langle1\rangle9$ of Remark 4.14 of [CGSS]; and it uses
2-subsets as the folded 5-cube does.

\begin{theorem}[SE2]\label{thm:se2}
Let $X$ be a strongly regular graph with parameters $(16,5,0,2)$. Then $X$ is isomorphic to the
folded 5-cube $\square_5$. Thus a strongly regular graph with parameters $(16,10,6,6)$ is
isomorphic to the Clebsch graph.
\end{theorem}

\begin{theorem}[SE3]\label{thm:se3}
Let $X$ be a strongly regular graph with parameters $(27,16,10,8)$. Then $X$ is isomorphic to the
Schl\"afli graph $\comp Q$. Equivalently, a strongly regular graph with parameters $(27,10,1,5)$
is isomorphic to $Q$.
\end{theorem}

\section{Complements}\label{sec:compl}

\begin{lemma}\label{lem:compl}
If $X$ is strongly regular with parameters $(n,k,a,c)$, then $\comp X$ is strongly regular with
parameters $(n,\ n-k-1,\ n-2-2k+c,\ n-2k+a)$, read in $\mathbb Z$.
\end{lemma}

A negative parameter here belongs to a condition that is vacuous: $n-2-2k+c$ is the number of
common neighbours in $\comp X$ of two vertices adjacent in $\comp X$, and $n-2k+a$ that of two
distinct vertices nonadjacent in $\comp X$, so either is negative only when there is no such pair,
as when $X$ is complete or empty.

\begin{proof}
Each vertex has $n-1-k$ neighbours in $\comp X$. Let $u\ne w$. A common neighbour of $u$ and $w$
in $\comp X$ is a vertex outside $\{u,w\}\cup\nb(u)\cup\nb(w)$. If $u\sim w$ in $X$, then $u,w\in
\nb(u)\cup\nb(w)$ and $\lvert\nb(u)\cup\nb(w)\rvert=2k-a$, so $u$ and $w$, which are
nonadjacent in $\comp X$, have $n-2k+a$ common neighbours there. If $u\not\sim w$ in $X$, then
$u,w\notin\nb(u)\cup\nb(w)$ and $\lvert\nb(u)\cup\nb(w)\rvert=2k-c$, so $u$ and $w$, which are
adjacent in $\comp X$, have $n-2-2k+c$ common neighbours there.
\end{proof}

This is \ml{SimpleGraph.IsSRGWith.compl}, in Mathlib's file
\ml{Mathlib/Combinatorics/SimpleGraph/StronglyRegular.lean}, which states the parameters as
\[
\text{\mlu{n (n - k - 1) (n - (2 * k - μ) - 2) (n - (2 * k - ℓ))}}
\]
in truncated subtraction. The two forms agree whenever the condition is not vacuous, and for the
four parameter sets below no subtraction is truncated (the smallest intermediate value is
$n-(2k-c)-2=0$, for $(16,10,6,6)$). The lemma turns $(16,5,0,2)$ and
$(16,10,6,6)$ into each other, and $(27,16,10,8)$ and $(27,10,1,5)$ into each other. Since
$\comp{\comp X}=X$, and a bijection that preserves adjacency and nonadjacency of distinct
vertices of $X$ and $Y$ does the same for $\comp X$ and $\comp Y$, we have $X\cong Y$ if and only
if $\comp X\cong\comp Y$.

\section{The folded 5-cube}\label{sec:se2}

\begin{lemma}\label{lem:fold-param}
The folded 5-cube is strongly regular with parameters $(16,5,0,2)$, and the Clebsch graph with
parameters $(16,10,6,6)$.
\end{lemma}

\begin{proof}
For $p\in V_{16}$ the map $s\mapsto s\sd p$ is a permutation of $V_{16}$ with $(s\sd p)\sd(t\sd
p)=s\sd t$, and so an automorphism of $\square_5$; it sends $p$ to $\emptyset$. Thus it suffices to
count at the vertex $\emptyset$. Its neighbours are the five 4-subsets $\Omega_5\diff\{i\}$. Two of
them have symmetric difference of size 2, so they are not adjacent, and every edge at $\emptyset$
lies in no triangle; thus $a=0$. A vertex not adjacent to $\emptyset$ and distinct from it is a
2-subset $P$, and $\lvert(\Omega_5\diff\{i\})\sd P\rvert$ is $4$ if $i\in P$ and $2$ if
$i\notin P$; so $\emptyset$ and $P$ have exactly the two common neighbours $\Omega_5\diff\{i\}$,
$i\in P$, and $c=2$. The second assertion follows by Lemma~\ref{lem:compl}.
\end{proof}

\begin{proof}[Proof of Theorem~\ref{thm:se2}]
Let $X$ be strongly regular with parameters $(16,5,0,2)$, fix a vertex $u$, and put $\nb=\nb(u)$
and $R=V(X)\diff(\{u\}\cup\nb)$. Then $\lvert\nb\rvert=5$ and $\lvert R\rvert=10$.

\emph{Step 1: $\nb$ is independent.} If $x,y\in\nb$ were adjacent, then $u$ would be a common
neighbour of the adjacent vertices $x$ and $y$, whereas $a=0$.

\emph{Step 2: the pairs.} Let $w\in R$. Then $w\ne u$ and $w\not\sim u$, so $w$ and $u$ have
exactly $c=2$ common neighbours, which are the neighbours of $w$ in $\nb$. We will write
$\varphi(w)=\nb(w)\cap\nb\in\binom{\nb}{2}$.

\emph{Step 3: $\varphi$ is a bijection $R\to\binom{\nb}{2}$.} Suppose that $w\ne w'$ in $R$ and
$\varphi(w)=\varphi(w')=\{x,y\}$. Then $x\ne y$, and $x\not\sim y$ by Step 1, while $u$, $w$ and
$w'$ are three distinct common neighbours of $x$ and $y$; this contradicts $c=2$. Thus $\varphi$ is
injective, and since $\lvert R\rvert=10=\binom52=\bigl\lvert\binom{\nb}{2}\bigr\rvert$, it is a
bijection.

\emph{Step 4: adjacency in $R$ is disjointness.} Let $w\in R$, with $\varphi(w)=\{x,y\}$. The five
neighbours of $w$ are $x$, $y$ and three vertices of $R$, since $w\not\sim u$. Let $w'\in R$ be a
neighbour of $w$. If $x\in\varphi(w')$, then $x$ would be a common neighbour of the adjacent
vertices $w$ and $w'$, whereas $a=0$; thus $x\notin\varphi(w')$, and likewise
$y\notin\varphi(w')$. Thus $\varphi(w')$ is one of the three 2-subsets of the 3-set
$\nb\diff\{x,y\}$. By Step 3, the three neighbours of $w$ in $R$ have three distinct such pairs,
so they are exactly the three vertices $w'\in R$ with $\varphi(w')\cap\varphi(w)=\emptyset$. Thus
for $w\ne w'$ in $R$ we have $w\sim w'$ if and only if $\varphi(w)\cap\varphi(w')=\emptyset$.

\emph{Step 5: the isomorphism.} Choose a bijection $\beta\colon\nb\to\Omega_5$, and define
$f\colon V(X)\to V_{16}$ by
\[
f(u)=\emptyset,\qquad f(x)=\Omega_5\diff\{\beta(x)\}\ (x\in\nb),\qquad
f(w)=\beta(\varphi(w))\ (w\in R).
\]
By Step 3, $f$ is a bijection onto $\{\emptyset\}\cup\{\text{4-subsets}\}\cup\binom{\Omega_5}{2}=
V_{16}$. We compare adjacency for distinct vertices $p,q$, by the parts that contain them.
\begin{enumerate}
\item $u$ and $x\in\nb$: adjacent in $X$, and $\lvert\emptyset\sd(\Omega_5\diff\{\beta(x)\})\rvert=
4$.
\item $u$ and $w\in R$: not adjacent, and $\lvert\emptyset\sd f(w)\rvert=2$.
\item $x,y\in\nb$: not adjacent by Step 1, and $\lvert f(x)\sd f(y)\rvert=\lvert\{\beta(x),
\beta(y)\}\rvert=2$.
\item $x\in\nb$ and $w\in R$: $x\sim w$ if and only if $x\in\varphi(w)$, that is, $\beta(x)\in
f(w)$; and, as in the proof of Lemma~\ref{lem:fold-param}, $\lvert f(x)\sd f(w)\rvert=4$ if and
only if $\beta(x)\in f(w)$.
\item $w,w'\in R$: $w\sim w'$ if and only if $f(w)\cap f(w')=\emptyset$ by Step 4, and for two
distinct 2-subsets $\lvert f(w)\sd f(w')\rvert=4-2\lvert f(w)\cap f(w')\rvert$ is 4 exactly when
they are disjoint.
\end{enumerate}
Thus $f$ is an isomorphism $X\to\square_5$.

For the consequence, let $Y$ be strongly regular with parameters $(16,10,6,6)$. By
Lemma~\ref{lem:compl}, $\comp Y$ is strongly regular with parameters $(16,5,0,2)$, so
$\comp Y\cong\square_5$, and $Y\cong\comp{\square_5}$ by Section~\ref{sec:compl}.
\end{proof}

\begin{remark}\label{rem:se2}
The proof uses $c=2$ twice (Steps 2 and 3), $a=0$ twice (Steps 1 and 4), and the valency 5 and
$n=16$ in Steps 3 and 4 (through $\lvert\nb\rvert=5$, $\lvert R\rvert=10$ and the three neighbours
of $w$ in $R$). The remaining instances of $c=2$ are then consequences. For $w\in R$ and
$z\in\nb\diff\varphi(w)$, the common neighbours of $w$ and $z$ are the two vertices $w'\in R$ with
$\varphi(w')=\{z,z'\}$, where $\{z,z'\}\subseteq\nb\diff\varphi(w)$. For distinct nonadjacent
$w,w'\in R$, whose pairs meet by Step 4 in one vertex $x$, the common neighbours of $w$ and $w'$
are $x$ and the vertex $w''\in R$ whose pair is $\nb\diff(\varphi(w)\cup\varphi(w'))$. Thus the
argument does not need them. The subgraph induced on $R$ is the Petersen graph (the Kneser graph of the 2-subsets of a
5-set), in agreement with [BvM] \S10.7; we do not use SE1.
\end{remark}

\section{The Schl\"afli graph}\label{sec:se3}

We will use the predicate of the blueprint (\ml{blueprint/sec2-lines.tex}, Definition \emph{the
triangle property}, label \ml{def:triangle-property}): a graph $Y$ has the \emph{triangle
property} if every edge $xy$ has exactly one common neighbour $z$, and every vertex outside
$\{x,y,z\}$ is adjacent to exactly one of $x,y,z$.

\begin{lemma}\label{lem:tri}
Let $Y$ be strongly regular with parameters $(n,k,1,c)$, where $n+3=3k$. Then $Y$ has the triangle
property.
\end{lemma}

The condition $n+3=3k$ is $n=3k-3$, and $n-3=3(k-2)$ as in the blueprint (Remark 4.14,
\ml{lem:414-GQ1}), written without subtraction, and so in the form of the Lean hypothesis
(Section~\ref{sec:lean}).

\begin{proof}
Let $x\sim y$. Since $a=1$, the edge $xy$ has exactly one common neighbour $z$, and $T=\{x,y,z\}$
is a triangle. A vertex $w\notin T$ is adjacent to at most one vertex of $T$: if it were
adjacent to two, say to $x$ and $y$ (the other cases are the same, since $xz$ and $yz$ are also
edges whose unique common neighbour is the third vertex of $T$), then $w$ and $z$ would be two
distinct common neighbours of $x$ and $y$. Each vertex of $T$ has exactly two neighbours in $T$,
and so $k-2$ neighbours outside $T$. Thus there are $3(k-2)=n-3$ edges between $T$ and
$V(Y)\diff T$, and no two of them end at the same vertex of $V(Y)\diff T$. Since $\lvert V(Y)\diff
T\rvert=n-3$, every vertex outside $T$ is the end of exactly one of them, that is, adjacent to
exactly one vertex of $T$.
\end{proof}

The condition $n+3=3k$ holds for $(9,4,1,2)$, $(15,6,1,3)$ and $(27,10,1,5)$, and more generally
for the collinearity graph of a generalized quadrangle of order $(2,t)$, with $n=3(2t+1)$ and
$k=2(t+1)$; the lemma is the converse half of GQ1. The parameter $c$ is not used.

\begin{lemma}\label{lem:q-param}
The graph $Q$ is strongly regular with parameters $(27,10,1,5)$, and the Schl\"afli graph with
parameters $(27,16,10,8)$.
\end{lemma}

\begin{proof}
The symmetric group on $\Omega_6$, acting on the indices, and the exchange $a_i\leftrightarrow
b_i$ are automorphisms of $Q$. Under the group they generate, every vertex lies in the orbit of
$a_1$ or of $c_{12}$, every edge in the orbit of $a_1b_2$, $a_1c_{12}$ or $c_{12}c_{34}$, and every
pair of distinct nonadjacent vertices in the orbit of $a_1a_2$, $a_1b_1$, $a_1c_{23}$ or
$c_{12}c_{13}$. The vertex $a_1$ has the ten neighbours $b_2,\dots,b_6$ and $c_{12},\dots,c_{16}$,
and $c_{12}$ has the ten neighbours $a_1,a_2,b_1,b_2$ and the six $c_{kl}$ with $\{k,l\}\subseteq
\{3,4,5,6\}$. The common neighbours are:
\begin{center}
\begin{tabular}{llc}
\toprule
pair & common neighbours & number\\
\midrule
$a_1b_2$ & $c_{12}$ & 1\\
$a_1c_{12}$ & $b_2$ & 1\\
$c_{12}c_{34}$ & $c_{56}$ & 1\\
\midrule
$a_1a_2$ & $b_3,b_4,b_5,b_6$, $c_{12}$ & 5\\
$a_1b_1$ & $c_{12},\dots,c_{16}$ & 5\\
$a_1c_{23}$ & $b_2$, $b_3$, $c_{14},c_{15},c_{16}$ & 5\\
$c_{12}c_{13}$ & $a_1$, $b_1$, $c_{45},c_{46},c_{56}$ & 5\\
\bottomrule
\end{tabular}
\end{center}
The second assertion follows by Lemma~\ref{lem:compl}.
\end{proof}

We will use the following result of the blueprint, which the repaired Lemma 2.10 of [CGSS]
applies in its step \ml{210unique}.

\begin{proposition}[{\ml{blueprint/sec2-lines.tex}, Lemma \emph{isomorphism from the order},
label \ml{lem:triangle-iso}}]\label{prop:triso}
Let $Y$ and $Z$ be finite graphs with the triangle property, each with an edge, neither with a
vertex adjacent to every other vertex, and with $\lvert V(Y)\rvert=\lvert V(Z)\rvert$. Then $Y\cong
Z$.
\end{proposition}

This is the graph-theoretic uniqueness assertion behind Lemma 2.10 of [CGSS], stated without a vector
representation; its Lean name is
\[
\text{\ml{SimpleGraph.HasTriangleProperty.nonempty_iso_of_card_eq}}
\]
(\ml{cgss/notation.md} \S10, row Lemma 2.10). We do not use the form
\[
\text{\mlu{LineSystem.Hyp24.orthGraph\_iso\_of\_card\_Γa\_eq}}
\]
of the uniqueness assertion, which asks that both graphs be orthogonality graphs of instances of Hypothesis 2.4; nor do we use the classification of
the generalized quadrangles with lines of size 3 (GQ2).

\begin{proof}[Proof of Theorem~\ref{thm:se3}]
Let $X$ be strongly regular with parameters $(27,16,10,8)$. By Lemma~\ref{lem:compl}, $\comp X$
is strongly regular with parameters $(27,10,1,5)$, and so is $Q$ by Lemma~\ref{lem:q-param}. Since
$27+3=3\cdot10$, both have the triangle property by Lemma~\ref{lem:tri}. Both have an edge, since
they are 10-regular on a nonempty set, and neither has a vertex adjacent to every other vertex,
since $10<26$. They have 27 vertices each. By Proposition~\ref{prop:triso}, $\comp X\cong Q$, and
so $X\cong\comp Q$ by Section~\ref{sec:compl}. The equivalent form is the same argument without
the first complementation.
\end{proof}

Thus SE3 in Lean is a reduction to the result of Section 2 of the blueprint: it needs Mathlib's
\ml{SimpleGraph.IsSRGWith.compl}, Lemma~\ref{lem:tri}, the parameters of $Q$, and
\ml{SimpleGraph.HasTriangleProperty.nonempty_iso_of_card_eq}. Lemma~\ref{lem:tri} also shows that
the strongly regular graphs with parameters $(9,4,1,2)$ and $(15,6,1,3)$ are unique, by the same
argument with $L_2(3)$ and the complement of $T(6)$ in place of $Q$ (the orthogonality graphs of
$\mathrm E_6$ and $\mathrm E_7$ in Section 2 of the blueprint; \ml{gq21} and \ml{gq22} in
\ml{cgss/notation.md} \S6).

\section{The computation}\label{sec:sage}

The script \ml{cgss/proofs/clebsch/checks/clebsch_check.sage}, run with SageMath~10, does the
following.
\begin{enumerate}
\item It builds $\square_5$ on the even subsets of $\{0,\dots,4\}$ and checks that it is strongly
regular with parameters $(16,5,0,2)$, that the Clebsch graph has parameters $(16,10,6,6)$ and is
the graph of symmetric differences of size 2, that $\square_5$ is isomorphic to
\ml{graphs.ClebschGraph()}, to the antipodal quotient of \ml{graphs.CubeGraph(5)}, and that the
Clebsch graph is isomorphic to the halved 5-cube.
\item At each of the 16 vertices $u$ of \ml{graphs.ClebschGraph()} it checks Steps 1 to 4 of the
proof of Theorem~\ref{thm:se2} ($\nb(u)$ independent, every $w\in R$ with two neighbours in
$\nb(u)$, $\varphi$ injective, three neighbours in $R$, adjacency in $R$ equal to disjointness of
the pairs), and that the map $f$ of Step 5 is an isomorphism onto $\square_5$.
\item It builds $Q$ and $\comp Q$ from the case lists of Definition~\ref{def:schlafli}, checks
that the two lists describe complementary graphs, that the parameters are $(27,10,1,5)$ and
$(27,16,10,8)$, that $\comp Q$ is isomorphic to \ml{graphs.SchlaefliGraph()}, and that $\comp Q$
is the graph of the $\mathrm E_8$ vectors of [BCN] 3.11 (ii) at inner product 1, and that the
correspondence of Section~\ref{sec:statement} is an isomorphism from $\comp Q$ onto the graph of
the vectors of [CGSS] Example 4.6 at inner product 1.
\item It checks the formula of Lemma~\ref{lem:compl} against Mathlib's truncated one for the four
parameter sets; the table of the proof of Lemma~\ref{lem:q-param}, pair by pair; that the
permutations of the indices and the exchange $a_i\leftrightarrow b_i$ are automorphisms of $Q$,
and that under the group they generate the vertices, the edges and the pairs of distinct
nonadjacent vertices fall into 2, 3 and 4 orbits, with the representatives of that proof.
\item It checks that the graph of the 16 vectors of [CGSS] Example 4.8 at inner product 1 is the
subgraph of the graph of Example 4.6 induced on the neighbours of $\mathbf e_5+\mathbf e_6$, and
that the map of Section~\ref{sec:lean} is an isomorphism from it onto $\comp{\square_5}$, the
certificate for Computation C9.
\item It checks the triangle property for $Q$ and for the complement of
\ml{graphs.SchlaefliGraph()}, the identity $n+3=3k$ for the three parameter sets of
Lemma 2.10 of [CGSS], that $Q$ has 45 triangles with 5 through each vertex, and that its automorphism
group has order $51840=\lvert W(\mathrm E_6)\rvert$.
\end{enumerate}
All checks pass. The computation is an independent check and not a part of the proof, which is
complete without it.

\section{The shape of the Lean statements}\label{sec:lean}
\raggedright

We will follow \ml{cgss/notation.md} \S10: the namespace \ml{SimpleGraph}, lowerCamelCase for
definitions, snake\_case theorems naming their conclusion, theorems about strongly regular graphs
in the namespace \ml{SimpleGraph.IsSRGWith} so that dot notation works, and \ml{nonempty_iso}
at the start of a theorem that an isomorphism exists (rule 1). Thus the uniqueness theorems are
\ml{IsSRGWith.nonempty_iso_clebsch} and not \ml{iso_clebsch}, which by rule 1 would be the name of
a definition building the isomorphism.

The graph names are those of \S6, with one definition for each name. The graphs
\ml{SimpleGraph.foldedFiveCube} and \ml{SimpleGraph.gq24} are the combinatorial graphs $\square_5$
and $Q$ of Definitions~\ref{def:clebsch} and~\ref{def:schlafli}. The graphs
\ml{SimpleGraph.clebsch} and \ml{SimpleGraph.schlafli} remain the graphs of the integer vectors of
[CGSS] Examples 4.8 and 4.6, as in Computation C9 of the blueprint
(\ml{blueprint/sec4-leasteigenvalue.tex}, \ml{comp:C9}), which gives them the parameters
$(16,10,6,6)$ and $(27,16,10,8)$ (\ml{clebsch_isSRGWith}, \ml{schlafli_isSRGWith}). We do not
define \mlu{clebsch := foldedFiveCubeᶜ} or \mlu{schlafli := gq24ᶜ}. Instead C9 supplies, as
certificates checked by \ml{decide}, the explicit isomorphisms
\begin{verbatim}
def SimpleGraph.clebsch.foldedFiveCubeComplIso : clebsch ≃g foldedFiveCubeᶜ
def SimpleGraph.schlafli.gq24ComplIso : schlafli ≃g gq24ᶜ
\end{verbatim}
named as \mlu{LineSystem.E₆.latticeGraphIso} and \mlu{LineSystem.E₇.triangularGraphComplIso} of
Computation C4 are: a definition that builds an isomorphism, in the namespace of the graph it
starts from, named after the other graph (rule 1). The second is the inverse of the
correspondence $a_i\mapsto\mathbf y_i$, $b_i\mapsto\mathbf z_i$, $c_{ij}\mapsto\mathbf x_{ij}$ of
Section~\ref{sec:statement}. The first, found by the script of Section~\ref{sec:sage}, sends
\[
\mathbf e_i+\mathbf e_5\mapsto s_i,\quad \mathbf e_i+\mathbf e_6\mapsto s_i\sd\{1,2\},\quad
\mathbf c-\mathbf e_i-\mathbf e_7\mapsto s_i\sd\{1,3,4,5\},\quad
\mathbf c-\mathbf e_i-\mathbf e_8\mapsto s_i\sd\{2,3,4,5\}
\]
for $1\le i\le4$, where $s_1=\emptyset$, $s_2=\{3,4\}$, $s_3=\{3,5\}$ and $s_4=\{4,5\}$.
\begin{center}
\begin{tabular}{ll}
\toprule
object & Lean \\
\midrule
$\square_5$ & \ml{SimpleGraph.foldedFiveCube : SimpleGraph S}, where \\
& \quad \ml{S = {s : Finset (Fin 5) // Even s.card}}, \\
& \quad with \mlu{Adj s t := (s.1 ∆ t.1).card = 4} \\
$\comp{\square_5}$ & \mlu{foldedFiveCubeᶜ}, with \\
& \quad \ml{SimpleGraph.compl_foldedFiveCube_adj_iff} ($\lvert s\sd t\rvert=2$) \\
$Q$ & \mlu{SimpleGraph.gq24 : SimpleGraph (Fin 2 × Fin 6 ⊕ T)}, where \\
& \quad \ml{T = {s : Finset (Fin 6) // s.card = 2}} \\
$\comp Q$ & \mlu{gq24ᶜ}, with \ml{SimpleGraph.compl_gq24_adj_iff} \\
Clebsch graph (Ex.~4.8) & \ml{SimpleGraph.clebsch} (C9), with \\
& \quad \ml{clebsch.foldedFiveCubeComplIso} \\
Schl\"afli graph (Ex.~4.6) & \ml{SimpleGraph.schlafli} (C9), with \ml{schlafli.gq24ComplIso} \\
Lemma~\ref{lem:fold-param} & \ml{SimpleGraph.foldedFiveCube_isSRGWith}, and \ml{.compl} \\
Lemma~\ref{lem:q-param} & \ml{SimpleGraph.gq24_isSRGWith}, and \ml{.compl} for $\comp Q$ \\
complements & \mlu{SimpleGraph.Iso.compl (f : X ≃g Y) : Xᶜ ≃g Yᶜ} \\
Lemma~\ref{lem:tri} & \ml{SimpleGraph.IsSRGWith.hasTriangleProperty} \\
\bottomrule
\end{tabular}
\end{center}
The lemmas \ml{compl_foldedFiveCube_adj_iff} and \ml{compl_gq24_adj_iff} state the adjacencies of
the complements given after Definitions~\ref{def:clebsch} and~\ref{def:schlafli}; they replace the
earlier proposals \ml{clebsch_adj_iff} and \ml{schlafli_adj_iff}, which named the complements when
these were called \ml{clebsch} and \ml{schlafli}. In \ml{gq24} the vertex $(0,i)$ is $a_{i+1}$,
$(1,i)$ is $b_{i+1}$, and a 2-subset $\{i,j\}$ is $c_{i+1,j+1}$; adjacency is by cases:
$(e,i)\sim(e',j)$ iff $e\ne e'$ and $i\ne j$, $(e,i)\sim s$ iff $i\in s$, and $s\sim t$ iff
\ml{Disjoint s.1 t.1}. Both vertex types have decidable equality and adjacency, and their
cardinalities, 16 and 27, follow by \ml{decide} or \ml{simp}. The docstrings record the
parameters, the names of the other sources (\S6), and the descriptions of
Section~\ref{sec:statement}. With \ml{X : SimpleGraph V}, \ml{[Fintype V] [DecidableEq V]
[DecidableRel X.Adj]}, the theorems are
\begin{verbatim}
theorem SimpleGraph.IsSRGWith.nonempty_iso_foldedFiveCube
    (h : X.IsSRGWith 16 5 0 2) : Nonempty (X ≃g foldedFiveCube)
theorem SimpleGraph.IsSRGWith.nonempty_iso_clebsch
    (h : X.IsSRGWith 16 10 6 6) : Nonempty (X ≃g clebsch)
theorem SimpleGraph.IsSRGWith.hasTriangleProperty
    (h : X.IsSRGWith n k 1 c) (hn : n + 3 = 3 * k) : X.HasTriangleProperty
theorem SimpleGraph.IsSRGWith.nonempty_iso_gq24
    (h : X.IsSRGWith 27 10 1 5) : Nonempty (X ≃g gq24)
theorem SimpleGraph.IsSRGWith.nonempty_iso_schlafli
    (h : X.IsSRGWith 27 16 10 8) : Nonempty (X ≃g schlafli)
\end{verbatim}
(the hypothesis \ml{hn} is $n+3=3k$ of Lemma~\ref{lem:tri}, without truncated subtraction). The
proof of \ml{nonempty_iso_foldedFiveCube} defines the map $f$ of Step 5 from a vertex \ml{u}, the
pair map \mlu{w ↦ X.neighborFinset u ∩ X.neighborFinset w} and an equivalence
\mlu{X.neighborFinset u ≃ Fin 5} (from the cardinality 5 and \ml{Finset.equivFinOfCardEq}); Steps
1 to 4 become four lemmas, and Step 5 a case split into the five cases. The proofs of
\ml{nonempty_iso_clebsch} and \ml{nonempty_iso_schlafli} apply \ml{IsSRGWith.compl}, then
\ml{nonempty_iso_foldedFiveCube} or \ml{nonempty_iso_gq24}, then \ml{Iso.compl} and
\ml{compl_compl}, which give \mlu{X ≃g foldedFiveCubeᶜ} or \mlu{X ≃g gq24ᶜ}, and compose with the
inverse of the certificate \ml{clebsch.foldedFiveCubeComplIso} or \ml{schlafli.gq24ComplIso}. The
proof of \ml{nonempty_iso_gq24} applies
\ml{SimpleGraph.HasTriangleProperty.nonempty_iso_of_card_eq} (Section 2 of the blueprint) to
\ml{X} and \ml{gq24}, with \ml{hasTriangleProperty} for both, the edge from the valency, and the
absence of a universal vertex (Mathlib's \ml{SimpleGraph.IsUniversal}) from $10<26$. The
parameters of the models may be proved by \ml{decide} (a check over at most $27^3$ triples), or by
the orbit argument of Lemmas~\ref{lem:fold-param} and~\ref{lem:q-param} if the kernel is too slow.

The Mathlib declarations that we use are in Mathlib at commit \ml{082e2d37} (4 September 2026),
the version pinned in \ml{lake-manifest.json}:
\begin{itemize}
\item \ml{Mathlib/Combinatorics/SimpleGraph/StronglyRegular.lean}: the structure
\ml{SimpleGraph.IsSRGWith} with the fields \ml{card}, \ml{regular}, \ml{of_adj}, \ml{of_not_adj};
\ml{SimpleGraph.IsSRGWith.compl}; \ml{IsSRGWith.card_neighborFinset_union_of_adj} and
\ml{IsSRGWith.card_neighborFinset_union_of_not_adj} (the counts of Lemma~\ref{lem:compl}).
\item \ml{Mathlib/Combinatorics/SimpleGraph/Basic.lean}: \ml{SimpleGraph.compl_adj},
\ml{SimpleGraph.commonNeighbors}, \ml{SimpleGraph.mem_commonNeighbors},
\ml{SimpleGraph.IsUniversal}; the instance \ml{SimpleGraph.completeAtomicBooleanAlgebra}, through
which the generic \ml{compl_compl} (\ml{Mathlib/Order/BooleanAlgebra/Basic.lean}) applies.
\item \ml{Mathlib/Combinatorics/SimpleGraph/Finite.lean}: \ml{SimpleGraph.neighborFinset},
\ml{SimpleGraph.card_neighborFinset_eq_degree}, \ml{SimpleGraph.IsRegularOfDegree},
\ml{IsRegularOfDegree.degree_eq}, \ml{IsRegularOfDegree.compl}.
\item \ml{Mathlib/Combinatorics/SimpleGraph/Maps.lean}: \ml{SimpleGraph.Iso} (notation \mlu{≃g}),
\ml{SimpleGraph.Iso.card_eq}, and \ml{SimpleGraph.Embedding.complEquiv}; Mathlib has the
complement of an embedding but no \ml{Iso.compl}, which is new (a few lines, from \ml{compl_adj}
and the injectivity of \ml{f}).
\item \ml{Mathlib/Order/SymmDiff.lean}: \ml{symmDiff}, with the scoped notation \mlu{∆}
(\ml{open symmDiff}); \ml{Mathlib/Data/Finset/Card.lean}: \ml{Finset.eq_of_subset_of_card_le};
\ml{Mathlib/Data/Fintype/EquivFin.lean}: \ml{Fintype.bijective_iff_injective_and_card} (Step 3);
\ml{Mathlib/Data/Finset/Powerset.lean}: \ml{Finset.card_powersetCard} (the count $\binom52=10$).
\item \ml{Mathlib/Combinatorics/SimpleGraph/Triangle/Basic.lean}: \ml{SimpleGraph.LocallyLinear},
which the triangle property implies (\S10, row Lemma 2.8); it is not needed here.
\end{itemize}

\section{Remarks on the sources}\label{sec:sources}

The direct argument for SE2 is Cameron's [C17, Theorem 2.2], and Seidel sketches a matrix proof
[S68, \S6, case 5]; we have written the direct argument in full. [BvM] \S10.7 and
\S10.10 state both uniqueness theorems without proof. [BCN] \S3.11 (iii) prints the parameters of the
complement of the Clebsch graph as $(16,5,0,1)$; they are
$(16,5,0,2)$, as \ml{cgss/prerequisites.md} (SE2) says. Where \ml{cgss/prerequisites.md} (SE3)
says that the graph is unique ``by GQ2'', we use instead the repaired uniqueness of Lemma 2.10, as
GQ3 allows: the triangle property of Lemma~\ref{lem:tri} is exactly what
Proposition~\ref{prop:triso} needs, and the lines of the quadrangle never appear.

\begin{thebibliography}{CGSS}
\bibitem[BCN]{BCN} A.~E. Brouwer, A.~M. Cohen and A. Neumaier, \emph{Distance-regular graphs},
Springer, 1989.
\bibitem[BvM]{BvM} A.~E. Brouwer and H. Van Maldeghem, \emph{Strongly regular graphs}, Cambridge
University Press, 2022.
\bibitem[C17]{C17} P.~J. Cameron, \emph{Eigenvalues and root systems in finite geometry}, lecture
notes, 2017.
\bibitem[CGSS]{CGSS} P.~J. Cameron, J.-M. Goethals, J.~J. Seidel and E.~E. Shult, Line graphs, root
systems, and elliptic geometry, \emph{J. Algebra} \textbf{43} (1976), 305--327.
\bibitem[S68]{S68} J.~J. Seidel, Strongly regular graphs with $(-1,1,0)$ adjacency matrix having
eigenvalue 3, \emph{Linear Algebra Appl.} \textbf{1} (1968), 281--298.
\end{thebibliography}

\end{document}
